\documentclass[10pt,a4paper]{article}
\usepackage[utf8]{inputenc}
\usepackage[portuguese]{babel}
\usepackage[margin=2cm]{geometry}
\usepackage{multicol}
\usepackage{enumitem}
\usepackage{xcolor}
\usepackage{listings}
\usepackage[box]{automultiplechoice}

% Configuração básica para exibição de códigos
\lstset{
  basicstyle=\small\ttfamily,
  breaklines=true
}

\begin{document}

{
\def\AMCbeginQuestion#1#2{\par\noindent{\bf (POSCOMP 2004 - Questão 26)}}
\begin{question}{} Em sistemas de memória virtual de paginação sob demanda, qual seria o critério ideal para substituição de páginas?
\begin{choices}
\wrongchoice{retirar a página que acabou de ser referenciada}
\correctchoice{retirar a página que será necessária no futuro mais distante}
\wrongchoice{retirar a página que está há mais tempo na memória}
\wrongchoice{retirar a página que foi referenciada menos vezes}
\wrongchoice{retirar a página que está há mais tempo sem ser utilizada} 
\end{choices}
\end{question}
}

{
\def\AMCbeginQuestion#1#2{\par\noindent{\bf (POSCOMP 2004 - Questão 27)}}
\begin{question}{} 27. Considere o seguinte programa com dois processos concorrentes. O escalonador poderá
alternar entre um e outro, isto é, eles poderão ser intercalados durante sua execução.
As variáveis x e y são compartilhadas pelos dois processos e inicializadas antes de sua
execução.
\begin{lstlisting}[language=C]
programa P
int x = 0;
int y = 0;
processo A {
  while (x == 0);
  print("a");
  y = 1;
  y = 0;
  print("d");
  y = 1;
}
processo B {
  print("b");
  x = 1;
  while (y == 0);
  print("c");
}
}
\end{lstlisting}
as possíveis saídas são:
\begin{choices}
\wrongchoice{adbc ou bcad}
\correctchoice{badc ou bacd}
\wrongchoice{abdc ou abcd}
\wrongchoice{dbca ou dcab}
\wrongchoice{Nenhuma das opções anteriores.} 
\end{choices}
\end{question}
}

{
\def\AMCbeginQuestion#1#2{\par\noindent{\bf (POSCOMP 2004 - Questão 39)}}
\begin{question}{}Em um sistema operacional, um processo pode, em um dado instante de tempo, estar
em um de três estados: em execução, pronto ou bloqueado. Considere as afirmativas abaixo sobre as possiveis transições entre estes estados que um processo pode realizar.
\begin{enumerate}[label=\Roman*.]
\item Do estado em \emph{execução} para o estado \emph{bloqueado}
\item Do estado em \emph{execução} para o estado \emph{pronto}
\item Do estado \emph{pronto} para o estado em \emph{execução}
\item Do estado \emph{pronto} para o estado \emph{bloqueado}
\item Do estado \emph{bloqueado} para o estado em \emph{execução}
\item Do estado \emph{bloqueado} para o estado \emph{pronto}
\end{enumerate}
\begin{choices}
\correctchoice{Somente as afirmativas I, II, III e VI são verdadeiras.}
\wrongchoice{Somente as afirmativas I, II e III são verdadeiras.}
\wrongchoice{Somente as afirmativas I, III, IV e VI são verdadeiras.}
\wrongchoice{Somente as afirmativas I, III, IV e V são verdadeiras.}
\wrongchoice{Todas as afirmativas são verdadeiras.} 
\end{choices}
\end{question}
}

{
\def\AMCbeginQuestion#1#2{\par\noindent{\bf (POSCOMP 2005 - Questão 36)}}
\begin{question}{} Quatro tarefas, A, B, C e D, estão prontas para serem executadas num único proces-
sador. Seus tempos de execução esperados são 9, 6, 3 e 5 segundos respectivamente. Em qual ordem eles devem ser executados para diminuir o tempo médio de resposta?
\begin{choices}
\correctchoice{C, D, B, A}
\wrongchoice{A, B, D, C}
\wrongchoice{C, B, D, A}
\wrongchoice{A, C, D, B}
\wrongchoice{O tempo médio de resposta independe da ordem.} 
\end{choices}
\end{question}
}

{
\def\AMCbeginQuestion#1#2{\par\noindent{\bf (POSCOMP 2005 - Questão 37)}}
\begin{question}{} Qual das alternativas a seguir melhor define uma Região Critica em Sistemas Operacionais?
\begin{choices}
\correctchoice{Um trecho de programa onde existe o compartilhamento de algum recurso que
não permite o acesso concomitante por mais de um programa.}
\wrongchoice{Um trecho de programa que deve ser executado em paralelo com a Região Critica
de outro programa.}
\wrongchoice{Um trecho de programa cujas instruções podem ser executadas em paralelo e em
qualquer ordem.}
\wrongchoice{Um trecho de programa onde existe algum recurso cujo acesso é dado por uma prioridade.}
\wrongchoice{Um trecho de
programa onde existe algum recurso a que somente o sistema operacional pode
ter acesso.} 
\end{choices}
\end{question}
}

{
\def\AMCbeginQuestion#1#2{\par\noindent{\bf (POSCOMP 2006 - Questão 32)}}
\begin{question}{}Qual dos seguintes mecanismos é o menos recomendado para se implementar
regiões criticas em sistemas operacionais?
\begin{choices}
\correctchoice{Espera ocupada}
\wrongchoice{Semáforo}
\wrongchoice{Troca de mensagens}
\wrongchoice{Monitores}
\wrongchoice{Variáveis de condição} 
\end{choices}
\end{question}
}


{
\def\AMCbeginQuestion#1#2{\par\noindent{\bf (POSCOMP 2006 - Questão 68)}}
\begin{question}{}A comunicação entre processos em um sistema distribuido pode ser realizada por
um mecanismo conhecido como RPC - chamada de procedimento remoto. Sobre este
mecanismo, assinale a opção correta abaixo:
\begin{choices}
\correctchoice{Os stubs cliente e servidor são responsáveis pela conversão de formato dos parâmetros
de entrada e saida, caso haja necessidade.}
\wrongchoice{Processos comunicantes compartilham o mesmo espaço de endereçamento.}
\wrongchoice{A geração dos stubs é comumente realizada por compilação a partir de uma es-
pecificação de interface realizada em uma linguagem de execução de interface
(IEL).}
\wrongchoice{O mecanismo faz uso de uma porta fixa, de número 8080, para comunicar difer-
entes processos e serviços entre computadores de um sistema distribuido.}
\wrongchoice{A falha de um cliente RPC gera uma chamada dita orfã no servidor que neste caso
repassa sempre os resultados do procedimento remoto para um proxy de retorno
especificado na chamada} 
\end{choices}
\end{question}
}


{
\def\AMCbeginQuestion#1#2{\par\noindent{\bf (POSCOMP 2006 - Questão 69)}}
\begin{question}{}Sobre algoritmos de exclusão mútua em sistemas distribuidos é correto afirmar
que:
\begin{choices}
\correctchoice{A maioria simples de permissões dos participantes para entrada em região critica
é suficiente para garantir a exclusão mútua no algoritmo distribuido.}
\wrongchoice{O algoritmo centralizado tem como principal desvantagem o alto número de troca
de mensagens.}
\wrongchoice{O algoritmo distribuido é totalmente independente da ordem dos eventos do sis-
tema distribuido.}
\wrongchoice{No algoritmo do token , a exclusão mútua é garantida por uma concessão de
bloqueio fornecida pelo gerente que mantém uma lista de tokens.}
\wrongchoice{Três mensagens são suficientes para fechar o ciclo de concessão, liberação e nova
concessão de acesso no algoritmo do token.} 
\end{choices}
\end{question}
}


{
\def\AMCbeginQuestion#1#2{\par\noindent{\bf (POSCOMP 2006 - Questão 70)}}
\begin{question}{}Um sistema distribuido pode manter diferentes cópias de um mesmo item de dado
a fim de melhorar o desempenho de leitura e aumentar a disponibilidade de acesso. A
modificação deste item de dado é realizada de acordo com protocolos de consistência
de cópias. Assinale a alternativa correta sobre esses protocolos.
\begin{choices}
\correctchoice{No protocolo de replicação ativa, todas as réplicas são atualizadas através de uma
única operação de escrita realizada por um mecanismo de multicast totalmente
ordenado.}
\wrongchoice{O protocolo baseado em cópia primária permite sempre a atualização da cópia
mais próxima e difunde o novo valor via unicast para todos os nós que mantém
uma outra cópia.}
\wrongchoice{A atualização de todas as cópias, no protocolo baseado em cópia primária, é
realizada através de um processo sincrono, onde o cliente é liberado para continuar
o fluxo de execução imediatamente após ter solicitado a atualização da cópia
primária.}
\wrongchoice{Nos protocolos baseados em quorum, os conflitos leitura-escrita e escrita-escrita
são evitados por autorizações de bloqueio (lock) emitidas por um coordenador
central ou sequenciador.}
\wrongchoice{Protocolos baseados em coerência de cache são mecanismos de consistência de
cópias que repassam a responsabilidade de manter essa consistência para os servi-
dores que detém cópias.} 
\end{choices}
\end{question}
}

{
\def\AMCbeginQuestion#1#2{\par\noindent{\bf (POSCOMP 2007 - Questão 35)}}
\begin{question}{}Analise as seguintes afirmativas e assinale a alternativa INCORRETA.
\begin{choices}
\correctchoice{O escalonamento de processos por prioridades utiliza múltiplas filas e garante que
todos os processos recebam sua fatia de tempo.}
\wrongchoice{O acesso a setores localizados em seqüência em uma mesma trilha de um disco
é mais rápido do que acessar o mesmo número de setores em trilhas diferentes,
devido ao menor número tanto de deslocamentos do cabeçote quanto de rotações
no disco.}
\wrongchoice{Na paginação por demanda, não é necessário que o processo inteiro se encontre
em memória para execução.}
\wrongchoice{O escalonamento de operações de entrada e saida em um disco rigido pode ser
utilizado para aumentar o desempenho. Porém, algoritmos como o SSTF (Shortest
Seek Time First) podem fazer com que requisições esperem indefinidamente.}
\wrongchoice{O surgimento do conceito de interrupções, juntamente com dispositivos de acesso
não-seqüencial, foi primordial para a evolução que levou aos sistemas multipro-
gramados.} 
\end{choices}
\end{question}
}

{
\def\AMCbeginQuestion#1#2{\par\noindent{\bf (POSCOMP 2007 - Questão 68)}}
\begin{question}{}Analise as seguintes afirmativas e assinale a alternativa INCORRETA.
\begin{enumerate}[label=\Roman*.]
\item Um “Servidor de Arquivos com Estado”, em um SAD, mantém todo seu estado
no caso de uma falha, garantindo a recuperação do mesmo sem a necessidade de
diálogo com os clientes.
\item II. Na gerência de cache em um SAD, uma das politicas utilizadas é a write-
through. O inconveniente dessa politica, comparada com outras, é a pouca confi-
abilidade no caso de falhas no cliente.
\item O uso de replicação em um SAD ao mesmo tempo que provê aumento na confia-
bilidade, também introduz um gargalo em termos de desempenho.
\end{enumerate}
A esse respeito, pode-se afirmar que
\begin{choices}
\correctchoice{nenhuma das afirmativas está correta.}
\wrongchoice{somente a afirmativa I está correta.}
\wrongchoice{somente a afirmativa II está correta.}
\wrongchoice{somente a afirmativa III está correta.}
\wrongchoice{somente as afirmativas I e III estão corretas.} 
\end{choices}
\end{question}
}

{
\def\AMCbeginQuestion#1#2{\par\noindent{\bf (POSCOMP 2007 - Questão 69)}}
\begin{question}{}Analise as seguintes afirmativas concernentes a questões de projeto de sistemas
distribuidos.
\begin{enumerate}[label=\Roman*.]
\item Um sistema distribuido tolerante a falhas deve continuar operando na presença
de problemas, podendo ocorrer uma degradação tanto no seu desempenho, como
nas suas funcionalidades.
\item No que diz respeito à escalabilidade, o projeto de um sistema distribuido deve
prever que a demanda nos serviços em qualquer dos equipamentos seja limitada
por uma constante dependente do número de nodos envolvidos.
\item Em um sistema distribuido transparente quanto à concorrência, a informação de
quantos usuários estão empregando determinado serviço deve ser omitida.
\end{enumerate}
A análise permite concluir que
\begin{choices}
\correctchoice{somente a afirmativa II está incorreta.}
\wrongchoice{somente a afirmativa I está incorreta.}
\wrongchoice{somente a afirmativa III está incorreta.}
\wrongchoice{somente as afirmativas I e III estão incorretas.}
\wrongchoice{todas as afirmativas estão incorretas.} 
\end{choices}
\end{question}
}

{
\def\AMCbeginQuestion#1#2{\par\noindent{\bf (POSCOMP 2007 - Questão 70)}}
\begin{question}{}Em relação aos sistemas distribuidos, analise as seguintes afirmativas.
\begin{enumerate}[label=\Roman*.]
\item Um sistema assincrono apresenta medida de tempo global.
\item A passagem de mensagens é o instrumento empregado para efetuar a comunica-
ção entre os processos de um sistema assincrono.
\item É possivel simular um computador paralelo de memória compartilhada usando-se
um sistema distribuido.
\item Quando um determinado elemento de um sistema distribuido efetua a difusão
de uma mensagem por meio de um multicast, todos os elementos do sistema
distribuido recebem a mensagem.
\end{enumerate}
A análise permite concluir que
\begin{choices}
\correctchoice{somente as afirmativas II e III estão corretas.}
\wrongchoice{somente a afirmativa IV está correta.}
\wrongchoice{somente as afirmativas I e II estão corretas.}
\wrongchoice{somente as afirmativas I e III estão corretas.}
\wrongchoice{somente as afirmativas I e IV estão corretas.} 
\end{choices}
\end{question}
}

\end{document}